\documentclass[a4paper,10pt,twocolumn,landscape]{article}
	\usepackage[ansinew]{inputenc}
	\usepackage{ae}
	\usepackage[brazil]{babel}
	\usepackage{xcolor}
	\usepackage{calc}
	\usepackage{wallpaper}
	\usepackage[margin=2cm]{geometry}
	\usepackage{SIunits}
	\usepackage{amsmath}
	\usepackage{paralist}
	\usepackage{enumerate}
	\usepackage{ifthen}

	% Minhas definições
	\newcounter{exercicio}
\newenvironment{exercicio}{%
	\refstepcounter{exercicio}%
	\penalty-200
	\noindent\colorbox{blue!60!black}{\makebox[\columnwidth-\fboxsep*2][c]{\textbf{\color{white}Exercício~\theexercicio}}}\smallskip
}{\par\medskip}
\newcommand{\foreign}[1]{\textsl{#1}}
\newcommand{\digite}[1]{{\fontfamily{cmss}\fontseries{bx}\selectfont#1}}	
\newcommand{\cs}[1]{{\normalfont\textbackslash\color{blue!50!black}#1}}
\newcommand{\pkg}[1]{{\normalfont\sffamily\color{orange}#1}}
\newcommand{\env}[1]{{\normalfont\sffamily\color{green!50!black}#1}}
\let\comando=\cs
\let\package=\pkg
\let\ambiente=\env

	% Papel de parede
	\CenterWallPaper{1}{../figuras/leao-pensador-wattermark.png}

	% Configuração extra: espaço entre as colunas.
	\setlength{\columnsep}{2cm}
	
\begin{document}

	
	\newcommand{\contador}[1]{\textit{#1}}
	\newcommand{\pacote}[1]{\textcolor{orange}{#1}}
	\newcommand{\etc}{\textsl{et cetera}} % Exercício 1
	
	%\newcommand{\foreign}[1]{\textsl{#1}} % Exercício 2
	\newcommand{\cs}[1]{\texttt{\textbackslash #1}} % Exercício 3
	\renewcommand{\etc}{\foreign{et cetera}}
	\renewcommand{\cs}[2][\null]{%
		\ifthenelse{\equal{#1}{\null}}%
		{\texttt{\textbackslash #2}}%
		{\texttt{\textbackslash #2\{#1\}}}%
	}

	\begin{center}\includegraphics[width=0.8\columnwidth]{../figuras/LogotipoCursoLaTeX_v2}\end{center}

	\section*{Comandos personalizados, listas e contadores}

	Monte um grupo de 2 ou 3 colegas e responda as questões abaixo. Envie o 
	arquivo \texttt{tex} com as respostas para o professor, informando os nomes
	dos integrantes do grupo como um comentário no arquivo de instruções.

	\bigskip

	% -----------
	% EXERCÍCIO 1
	% -----------
	\begin{exercicio}
	Crie o comando \cs{etc}, que imprime \foreign{et cetera} na fonte com forma
	inclinada (\foreign{slanted}), exatamente como aqui. Depois, teste o seu
	comando.
	\end{exercicio}
	
	% -----------
	% EXERCÍCIO 2
	% -----------
	\begin{exercicio}
	Às vezes é útil destacar levemente palavras que não fazem parte do idioma
	utilizado num texto. \foreign{C'est la vie}, por exemplo. Crie o comando
	\cs{foreign} (do inglês, ``estrangeiro''), que imprime seu argumento na
	fonte de forma inclinada (\foreign{slanted}).
	\end{exercicio}
		
	% -----------
	% EXERCÍCIO 3
	% -----------
	\begin{exercicio}
	\foreign{Et cetera} também não faz parte do português. Logo, é recomendado
	que o próprio comando \cs{etc} seja definido em termos de \cs{foreign}. Faça
	isso: \emph{redefina} o comando \cs{etc} utilizando \cs{foreign} em sua
	definição.
	\end{exercicio}
	
	% -----------
	% EXERCÍCIO 4
	% -----------
	\begin{exercicio}
	Embora a barra-invertida não seja muito utilizada, ainda assim é bastante 
	chato ficar escrevendo \cs{textbackslash}. Isto é
	especialmente verdadeiro num curso de \LaTeX, onde é preciso imprimir os
	comandos --- com a barra-invertida --- diversas vezes. Ajude-me, então: crie
	o comando \cs{cs} (de \foreign{control sequence}) que recebe como seu único
	argumento o nome do comando. Como resultado, este comando deve imprimir a
	barra-invertida e o nome do comando, utilizando as fontes da família	
	{\ttfamily Computer Modern Typewriter}. Por exemplo, o comando \cs[emph]{cs}
	deve imprimir \cs{emph}.
	\end{exercicio}
	
	% -----------
	% EXERCÍCIO 5
	% -----------
	\begin{exercicio}	
	\textbf{Re}defina o comando \cs{cs} para que ele possa representar os 
	comandos
	com um argumento obrigatório. Isto é, redefina-o de modo que seu primeiro 
	argumento seja o nome do comando e o segundo, o argumento dele. Por exemplo,
	o comando \verb|\cs{color}{red}| deve imprimir \cs[red]{color}.
	\end{exercicio}
	
	% -----------
	% EXERCÍCIO 6
	% -----------
	\begin{exercicio}
	Imagine-se como o autor de um livro sobre o \LaTeX, no qual você utiliza
	intensa --- e sabiamente --- o comando \cs{cs}. Mas ocorre que na metade do
	trabalho, por alguma razão você decide que convém escrever o \emph{nome} dos
	comandos em azul (a barra-invertida, chaves e argumentos devem permanecer em
	preto). Pois bem: redefina \cs{cs} para fazer isso (não se esqueça do pacote
	\pacote{color} ou \pacote{xcolor}).
	
	\renewcommand{\cs}[2][\null]{
		\ifthenelse{\equal{#1}{\null}}
		{\texttt{\textbackslash{\color{blue} #2}}}
		{\texttt{\textbackslash{\color{blue} #2}\{#1\}}}
	}
	
	eg: \verb|\cs{color}{red}| deve imprimir \cs[red]{color}.	
	\end{exercicio}
	%\renewcommand{\cs}[2]{\texttt{\textbackslash{\color{blue}#1}\{#2\}}}
	%\cs{begin}{enumerate}

	% -----------
	% EXERCÍCIO 7
	% -----------
	\begin{exercicio}
	Saber buscar informações sobre um pacote e aplicá-las é uma tarefa rotineira
	no \TeX, \LaTeX\ e outros. E é isso que você praticará nesta tarefa: procure 
	na Internet as informações necessárias para utilizar o \emph{pacote}
	\pacote{enumerate}, que agrega melhorias à \emph{lista} ``enumerate.'' 
	Utilize-o para montar a lista abaixo (atenção aos números que precedem cada
	item, que devem ser reproduzidos fielmente):
	
	\begin{enumerate}[(i)] % <------------------------
		\item Maçã é vermelha
		\item Uva é roxa
		\item Banana é amarela
		\item Laranja é\dots\ laranja!
	\end{enumerate}	
	\end{exercicio}

	% -----------
	% EXERCÍCIO 8
	%
	% Criando sua própria lista.
	% --------------------------
	\begin{exercicio}
	Para terminar, crie uma lista de nome \texttt{ENUMERATE} que produz 	
	exatamente a mesma lista do exercício anterior, mas sem utilizar o pacote
	\pacote{enumerate}. Ou seja, defina uma lista genérica (ambiente ``list'') 
	com o comando \cs{newenvironment} cujos números que precedem os ítens são
	(i), (ii), \etc.
	
	\textbf{obs.:} você não precisa criar um contador para esta lista. Ao invés
	disso, use o próprio \contador{enumi} (da lista ``enumerate''). Depois sim:
	tente criar seu próprio contador para esta lista.
	\end{exercicio}
	% --------
	% RESPOSTA
	% --------
%	\newenvironment{ENUMERATE}{%
%		\begin{list}{%
%			(\roman{ENUMI})% <----------------------------
%		}{
%			\newcounter{ENUMI}
%			\usecounter{ENUMI}
%		}
%	}{\end{list}}
%
%	\begin{ENUMERATE}
%		\item Maçã é vermelha
%		\item Uva é roxa
%		\item Banana é amarela
%		\item Laranja é\dots\ laranja!
%	\end{ENUMERATE}
	 
		
\end{document}
